\subsection*{CU-10: Gestionar sesiones activas}
\addcontentsline{toc}{subsection}{CU-10: Gestionar sesiones activas}
\label{cu-10}

\begin{fichacu}{CU-10}{Gestionar sesiones activas}
  \crudFila{Responsable}{José Eduardo Prior Hernández}
  \crudFila{Actor principal}{Jugador}
  \crudFila{Actores de sistema}{Servidor de partidas, Base de datos}
  \crudFila{Prototipos}{9f}
  \crudFila{Objetivo}{Mostrar al Jugador en qué dispositivos tiene la sesión abierta y permitirle cerrar a distancia una o todas las demás, por ejemplo cuando sospecha que alguien entró a su cuenta o dejó la sesión abierta en un equipo ajeno.}
  \crudFila{Disparador}{El Jugador selecciona ``sesiones activas'' en la \texttt{GUI\_\allowbreak{}AccountSettings}.}
  \textbf{Precondiciones} & \cuCelda{\begin{cureglas}
      \item \textbf{\mbox{PRE-01}} El Jugador tiene una \textit{Sesion} abierta y su token es válido.
      \item \textbf{\mbox{PRE-02}} El Servidor de partidas está en ejecución y tiene conexión disponible con la Base de datos.
    \end{cureglas}} \\ \hline
  \textbf{Flujo normal} & \cuCelda{\begin{cuenum}[start=1]
      \item El Jugador selecciona ``sesiones activas'' en la \texttt{GUI\_\allowbreak{}AccountSettings}.
      \item El SJM envía el mensaje \texttt{LIST\_\allowbreak{}SESSIONS\_\allowbreak{}REQUEST} con el token de sesión. \textit{(Ver \mbox{EX-02}, \mbox{EX-03}, \mbox{EX-04})}
      \item El Servidor de partidas verifica que el token corresponda a una \textit{Sesion} abierta. \textit{(Ver \mbox{FA-01})}
      \item El Servidor de partidas recupera las \textit{Sesion} del \textit{Usuario} con \texttt{fecha\_\allowbreak{}fin} nula y \texttt{fecha\_\allowbreak{}expiracion} posterior a la fecha actual (\mbox{RN-02}). \textit{(Ver \mbox{EX-01})}
      \item El Servidor de partidas identifica cuál de ellas corresponde al token de la solicitud.
      \item El Servidor de partidas responde con \texttt{LIST\_\allowbreak{}SESSIONS\_\allowbreak{}OK} y, por cada sesión, su identificador, \texttt{nombre\_\allowbreak{}dispositivo}, \texttt{direccion\_\allowbreak{}ip} parcialmente oculta, \texttt{fecha\_\allowbreak{}inicio}, \texttt{fecha\_\allowbreak{}ultimo\_\allowbreak{}uso} y la marca de sesión actual.
      \item El SJM despliega la \texttt{GUI\_\allowbreak{}ActiveSessions} con la lista ordenada por último uso, la sesión actual destacada y sin opción de cerrarse desde aquí, un botón ``cerrar'' en cada una de las demás y la opción ``cerrar todas las demás''. \textit{(Ver \mbox{FA-02})}
      \item El Jugador selecciona ``cerrar'' en una sesión. \textit{(Ver \mbox{FA-03}, \mbox{FA-04})}
    \end{cuenum}} \\
   & \cuCelda{\begin{cuenum}[start=9]
      \item El SJM despliega la \texttt{GUI\_\allowbreak{}MessageConfirm} con el nombre del dispositivo y la fecha de último uso, con las opciones ``Cerrar sesión'' y ``Cancelar''. \textit{(Ver \mbox{FA-05})}
      \item El Jugador confirma.
      \item El SJM envía el mensaje \texttt{CLOSE\_\allowbreak{}SESSION\_\allowbreak{}REQUEST} con el identificador de la sesión a cerrar y el token.
      \item El Servidor de partidas verifica que la sesión indicada pertenezca al mismo \textit{Usuario} del token (\mbox{RN-01}). \textit{(Ver \mbox{FA-06})}
      \item El Servidor de partidas verifica que la sesión siga abierta. \textit{(Ver \mbox{FA-07})}
      \item El Servidor de partidas registra en la \textit{Sesion} la \texttt{fecha\_\allowbreak{}fin} actual y el \texttt{motivo\_\allowbreak{}cierre} igual a \texttt{CIERRE\_\allowbreak{}REMOTO}, e invalida su token. \textit{(Ver \mbox{EX-01})}
      \item El Servidor de partidas registra el cierre en la \textit{BitacoraAcceso} con el resultado \texttt{CIERRE\_\allowbreak{}REMOTO}.
      \item El Servidor de partidas notifica al cliente de la sesión cerrada, si sigue conectado, con \texttt{SESSION\_\allowbreak{}CLOSED\_\allowbreak{}REMOTELY} (\mbox{RN-04}). \textit{(Ver \mbox{FA-08})}
    \end{cuenum}} \\
   & \cuCelda{\begin{cuenum}[start=17]
      \item El Servidor de partidas responde con \texttt{CLOSE\_\allowbreak{}SESSION\_\allowbreak{}OK}.
      \item El SJM retira la sesión de la lista y muestra la confirmación junto a ella.
      \item Fin del caso de uso.
    \end{cuenum}} \\ \hline
  \textbf{Flujos alternos} & \cuCelda{\textbf{\mbox{FA-01}: La sesión desde la que se consulta ya no es válida}\par
    \begin{cuenum}[start=1]
      \item El Servidor de partidas responde con \texttt{SESSION\_\allowbreak{}INVALID}: pudo cerrarse desde otro dispositivo mientras el Jugador tenía abierta la pantalla.
      \item El SJM muestra la \texttt{GUI\_\allowbreak{}MessageWarning} indicando que su sesión terminó y despliega la \texttt{GUI\_\allowbreak{}Login}.
      \item Fin del caso de uso.
    \end{cuenum}} \\
   & \cuCelda{\textbf{\mbox{FA-02}: Solo existe la sesión actual}\par
    \begin{cuenum}[start=1]
      \item El SJM muestra la sesión actual y un estado vacío indicando que no hay otros dispositivos con la sesión abierta.
      \item El SJM no ofrece la opción ``cerrar todas las demás''.
      \item Fin del caso de uso.
    \end{cuenum}} \\
   & \cuCelda{\textbf{\mbox{FA-03}: El Jugador cierra todas las demás sesiones}\par
    \begin{cuenum}[start=1]
      \item El Jugador selecciona ``cerrar todas las demás''.
      \item El SJM despliega la \texttt{GUI\_\allowbreak{}MessageConfirm} indicando cuántos dispositivos se desconectarán y que la sesión actual se conserva.
      \item El Jugador confirma y el SJM envía el mensaje \texttt{CLOSE\_\allowbreak{}OTHER\_\allowbreak{}SESSIONS\_\allowbreak{}REQUEST}.
      \item El Servidor de partidas cierra todas las \textit{Sesion} abiertas del \textit{Usuario} distintas de la del token, con \texttt{motivo\_\allowbreak{}cierre} igual a \texttt{CIERRE\_\allowbreak{}REMOTO}, en una sola transacción. \textit{(Ver \mbox{EX-01})}
      \item El Servidor de partidas registra un evento en la \textit{BitacoraAcceso} por cada sesión cerrada y notifica a los clientes conectados.
      \item El Servidor de partidas responde con \texttt{CLOSE\_\allowbreak{}OTHER\_\allowbreak{}SESSIONS\_\allowbreak{}OK} y el número de sesiones cerradas.
      \item El SJM deja en la lista únicamente la sesión actual y propone cambiar la contraseña, que continúa en el \mbox{CU-09} si el Jugador lo acepta (\mbox{RN-05}).
      \item Fin del caso de uso.
    \end{cuenum}} \\
   & \cuCelda{\textbf{\mbox{FA-04}: El Jugador quiere cerrar la sesión actual}\par
    \begin{cuenum}[start=1]
      \item El SJM no ofrece ``cerrar'' sobre la sesión actual, e indica junto a ella que para salir de este dispositivo se usa ``cerrar sesión''.
      \item Si el Jugador selecciona esa opción, el flujo continúa en el \mbox{CU-05} Cerrar sesión.
      \item Fin del caso de uso.
    \end{cuenum}} \\
   & \cuCelda{\textbf{\mbox{FA-05}: El Jugador cancela el cierre}\par
    \begin{cuenum}[start=1]
      \item El Jugador selecciona ``Cancelar''.
      \item El SJM cierra la confirmación sin enviar nada.
      \item El flujo continúa en el paso 8 del FN.
    \end{cuenum}} \\
   & \cuCelda{\textbf{\mbox{FA-06}: La sesión indicada no pertenece al Jugador}\par
    \begin{cuenum}[start=1]
      \item El Servidor de partidas responde con \texttt{CLOSE\_\allowbreak{}SESSION\_\allowbreak{}FAILED} y el motivo \texttt{SESION\_\allowbreak{}NO\_\allowbreak{}ENCONTRADA}, sin revelar si existe una sesión con ese identificador en otra cuenta.
      \item El Servidor de partidas registra el intento en la bitácora del servidor con la \texttt{version\_\allowbreak{}cliente} y la dirección de red de origen, por tratarse de un indicio de cliente modificado.
      \item El SJM recarga la lista.
      \item El flujo continúa en el paso 7 del FN.
    \end{cuenum}} \\
   & \cuCelda{\textbf{\mbox{FA-07}: La sesión ya estaba cerrada o expiró}\par
    \begin{cuenum}[start=1]
      \item El Servidor de partidas no modifica la \textit{Sesion} y responde con \texttt{CLOSE\_\allowbreak{}SESSION\_\allowbreak{}OK}, porque el estado que el Jugador pedía ya se cumple.
      \item El flujo continúa en el paso 18 del FN.
    \end{cuenum}} \\
   & \cuCelda{\textbf{\mbox{FA-08}: El dispositivo cerrado estaba jugando una partida}\par
    \begin{cuenum}[start=1]
      \item El cliente de la sesión cerrada recibe \texttt{SESSION\_\allowbreak{}CLOSED\_\allowbreak{}REMOTELY}, abandona la pantalla de partida y despliega la \texttt{GUI\_\allowbreak{}Login} con el aviso de que la sesión se cerró desde otro dispositivo.
      \item La \textit{Participacion} del Jugador \textbf{no} se cierra y su reloj sigue corriendo: cerrar una sesión no es abandonar (\mbox{RN-06}).
      \item El SJM del dispositivo actual ofrece volver a esa partida, lo que continúa en el \mbox{CU-26}.
      \item El flujo continúa en el paso 17 del FN.
    \end{cuenum}} \\ \hline
  \textbf{Excepciones} & \cuCelda{\textbf{\mbox{EX-01}: Fallo al consultar o actualizar la base de datos}\par
    \begin{cuenum}[start=1]
      \item El Servidor de partidas revierte la transacción en curso, si la hubiera, de modo que en el \mbox{FA-03} no queden unas sesiones cerradas y otras no.
      \item El Servidor de partidas registra la excepción en la bitácora del servidor con un identificador de incidencia y responde con \texttt{SERVER\_\allowbreak{}ERROR}.
      \item El SJM muestra la \texttt{GUI\_\allowbreak{}MessageError} con el identificador.
      \item El flujo continúa en el paso 7 del FN, con la lista tal como estaba.
    \end{cuenum}} \\
   & \cuCelda{\textbf{\mbox{EX-02}: Se agota el tiempo de espera de la respuesta}\par
    \begin{cuenum}[start=1]
      \item El SJM detecta que transcurrió el tiempo límite sin recibir un mensaje completo.
      \item El SJM no reenvía la solicitud de cierre y ofrece recargar la lista, que dirá si el cierre se aplicó.
      \item El flujo continúa en el paso 2 del FN si el Jugador recarga.
    \end{cuenum}} \\
   & \cuCelda{\textbf{\mbox{EX-03}: La conexión se interrumpe durante el intercambio}\par
    \begin{cuenum}[start=1]
      \item El SJM muestra la barra de conexión perdida.
      \item Al recuperar la conexión, el flujo continúa en el paso 2 del FN.
    \end{cuenum}} \\
   & \cuCelda{\textbf{\mbox{EX-04}: El mensaje recibido está malformado}\par
    \begin{cuenum}[start=1]
      \item El SJM descarta el mensaje, cierra la conexión y registra en la bitácora local el contenido truncado.
      \item El SJM muestra la \texttt{GUI\_\allowbreak{}MessageError} indicando que la respuesta no pudo interpretarse.
      \item Fin del caso de uso.
    \end{cuenum}} \\ \hline
  \textbf{Postcondiciones} & \cuCelda{\begin{cureglas}
      \item \textbf{\mbox{POST-01}} Si el caso termina por el FN, la \textit{Sesion} elegida tiene \texttt{fecha\_\allowbreak{}fin} y \texttt{motivo\_\allowbreak{}cierre} igual a \texttt{CIERRE\_\allowbreak{}REMOTO}, y su token ya no es válido.
      \item \textbf{\mbox{POST-02}} Si el caso termina por el \mbox{FA-03}, la única \textit{Sesion} abierta de la cuenta es la actual.
      \item \textbf{\mbox{POST-03}} La sesión desde la que se opera nunca se cierra en este caso.
      \item \textbf{\mbox{POST-04}} Ninguna \textit{Participacion} cambió: una partida en curso sigue viva aunque se cierre la sesión del dispositivo que la jugaba.
      \item \textbf{\mbox{POST-05}} Cada sesión cerrada quedó registrada en la \textit{BitacoraAcceso}.
    \end{cureglas}} \\ \hline
  \textbf{Reglas de negocio} & \cuCelda{\begin{cureglas}
      \item \textbf{\mbox{RN-01}} Un jugador solo puede ver y cerrar sus propias sesiones.
      \item \textbf{\mbox{RN-02}} Solo se listan sesiones con \texttt{fecha\_\allowbreak{}fin} nula y \texttt{fecha\_\allowbreak{}expiracion} futura. Una sesión expirada no se muestra aunque nunca se haya cerrado explícitamente.
      \item \textbf{\mbox{RN-03}} La dirección de red se muestra parcialmente oculta: basta para reconocer un lugar, no para identificarlo.
      \item \textbf{\mbox{RN-04}} Cerrar una sesión a distancia invalida su token en el servidor. El dispositivo afectado lo descubre en su siguiente mensaje, o de inmediato si sigue conectado.
      \item \textbf{\mbox{RN-05}} Tras cerrar todas las demás sesiones se propone cambiar la contraseña, porque si alguien las abrió sin permiso, las volverá a abrir con la misma contraseña.
      \item \textbf{\mbox{RN-06}} Cerrar una sesión no retira al jugador de una partida en curso ni detiene su reloj (\mbox{RN-04} del \mbox{CU-05}).
      \item \textbf{\mbox{RN-07}} Este caso existe porque una cuenta puede tener varias sesiones abiertas a la vez (\mbox{D-01}).
    \end{cureglas}} \\ \hline
  \textbf{Observación} & \cuCelda{\textit{Sobre por qué se necesita fecha\_ultimo\_uso.}\par
    La \texttt{fecha\_\allowbreak{}inicio} de una sesión dice cuándo se abrió, pero no si todavía se usa: una sesión abierta hace veinte horas en el equipo de casa y una abierta hace veinte horas en un equipo ajeno se ven iguales. Mostrar la fecha de último uso es lo que permite al jugador reconocer una sesión olvidada. El servidor la actualiza al atender mensajes de esa sesión, con una granularidad de minutos para no escribir en cada jugada.} \\ \hline
  \textbf{Datos que toca} & \cuCelda{\begin{cudatos}
      \item \textit{Sesion}: lee las abiertas del usuario con \texttt{nombre\_\allowbreak{}dispositivo}, \texttt{direccion\_\allowbreak{}ip}, \texttt{fecha\_\allowbreak{}inicio}, \texttt{fecha\_\allowbreak{}ultimo\_\allowbreak{}uso}, \texttt{fecha\_\allowbreak{}expiracion} \textit{(paso 4)}
      \item \textit{Sesion}: escribe \texttt{fecha\_\allowbreak{}fin} y \texttt{motivo\_\allowbreak{}cierre} \textit{(pasos 14, \mbox{FA-03})}
      \item \textit{BitacoraAcceso}: inserta \textit{(pasos 15, \mbox{FA-03})}
      \item \textit{Participacion}: lee, para el \mbox{FA-08} \textit{(paso \mbox{FA-08})}
    \end{cudatos}} \\ \hline
\end{fichacu}
\clearpage
